\documentclass{article}
\usepackage[utf8]{inputenc}
\usepackage[T1]{fontenc}
\usepackage{geometry}
\usepackage{graphicx}
\usepackage{float}
\usepackage{hyperref}
\usepackage[english,indonesian]{babel}

\title{Optimasi Pemrosesan Template LaTeX Menggunakan Pendekatan Berbasis AI}
\author{Academic Support Agent}
\date{\today}

\begin{document}
\maketitle

\begin{abstract}

\end{abstract}

\chapter{PENDAHULUAN}


\chapter{TINJAUAN PUSTAKA}


% Original content (commented out):
% \documentclass[12pt,a4paper]{article}
%%%%%%%%%%%%%%%%%%%%%%%%%%%%%%%%%%%%%%%%%%%%%%%%%%%%%
%\usepackage[utf8]{inputenc}
\usepackage[portuguese]{babel}
\addto{\captionsportuguese}{%
  \renewcommand{\refname}{\hfill \large REFERÊNCIAS \hfill} % V V V
}
\addto{\captionsportuguese}{% 
  \renewcommand{\contentsname}{\hfill SUMÁRIO \hfill}} %\centering não tava funcionando sabe deus por que, usei o \hfill pra enjambrar
    
\usepackage[T1]{fontenc}
\usepackage{fontspec}
\usepackage{amsmath,esint}
\usepackage{amsfonts}
\usepackage{amssymb}
\usepackage{makeidx}
\usepackage{graphicx}
\usepackage{lettrine}
\usepackage{lmodern}
\usepackage{dirtytalk}
\usepackage{setspace}
    \linespread{1.5}
\usepackage{indentfirst}
    \setlength{\parindent}{1cm}
    \setlength{\parskip}{0pt}
\usepackage{pdfpages}
\usepackage[left=3cm,right=2cm,top=3cm,bottom=2cm]{geometry}
\usepackage{hyperref}
\usepackage{multirow}
\usepackage{booktabs}
\usepackage{subfigure}
\usepackage{fancyhdr}
\usepackage{soul} % Pacote que permite grifar palavras no texto

\usepackage[ %faz a citação ser do tipo autor data, faz parecer et al na citação, mas na referencia mostra todos os autores e deixa o titulo em negrito na referencia
	alf,
	abnt-repeated-title-omit=yes,	
	abnt-emphasize=bf,
	abnt-etal-list=0
	]
	{abntex2cite}
\citebrackets()%faz a citação ser entre parenteses e não entre colchetes

%-----------------------------------------------------------
\hypersetup{
%links coloridos
   pdfborder={0 0 0},
   colorlinks = true,  
   linkcolor = black,          % Cor dos links internos
   citecolor = blue        % Cor dos links das referências bibliográficas
}
\setmainfont{Times New Roman}
\pagestyle{fancy}
% Clear the header and footer
\fancyhead{}
\fancyfoot{}
\fancyhead[R]{\thepage}
\renewcommand{\headrulewidth}{0pt}%
\setlength{\headheight}{14.49998pt}
%-----------------------------------------------------------

%------Redefinido Tamanho de Fontes-------------------------

\renewcommand{\large}{\fontsize{14}{15}\selectfont}  % large vira 14pt
\usepackage{titlesec}
    \titleformat{\section}
      {\normalfont\fontsize{14}{15}\bfseries}{\thesection}{1em}{}

%-----------------------------------------------------------


\begin{document}

%___________________________CAPA___________________________%
    \begin{titlepage}

        \begin{center}
                \textbf{INSTITUTO FEDERAL CATARINENSE - RIO DO SUL}
            
                \textbf{CURSO DE LICENCIATURA EM FÍSICA} 
        \end{center}
   
    \vspace*{2cm}
    \begin{center}

        \textbf{NOME}
       
    \end{center}
    \vspace*{5cm}

    \begin{center}
    \large
         \textbf{TITULO DO RELATÓRIO}
    \end{center}
      
    \vfill  
    
    \begin{center}
        \textbf{Rio do Sul}

        \textbf{ano}
        
    \end{center}
    
    

    \end{titlepage}



%______________________Folha de Rosto__________________________%

\thispagestyle{empty} %sem numeração de página

    \begin{center}
        \textbf{NOME}
    \end{center}

    \vspace{3 cm}

    \begin{center}
        \large
        \textbf{TÍTULO DO RELATÓRIO DE ESTÁGIO}
    \end{center}

    \vspace{3 cm}
    
    \hfill\begin{minipage}{\dimexpr\textwidth-8,5 cm}
        \linespread{1,0}\selectfont
        \large 
            Relatório de Estágio Supervisionado I do Curso de Licenciatura em Física do Instituto Federal Catarinense Campus Rio do Sul realizado na Escola XXXXXXXXXXX,  como requisito parcial para conclusão da disciplina de Estágio Supervisionado I.\\

            Professor Orientador: XXX

    \xdef\tpd{\the\prevdepth}
    \end{minipage}

    \vfill

    \begin{center}
        \textbf{Rio do Sul\\
        Ano}
    \end{center}





\newpage

%______________________Folha em Branco__________________________%
\thispagestyle{empty} 

\

\newpage

%______________________________Sumário_________________________%
\thispagestyle{empty} 
\tableofcontents
\newpage
%_____________________________Introdução__________________________%
\section{INTRODUÇÃO}

        \textcolor{red}{Deve iniciar com um parágrafo abrangente. Por exemplo, escrever sobre o ensino médio, ensino de Física, a escola como espaço formal de educação, etc.}
        
        \textcolor{red}{É importante manifestar concepções sobre educação, escola, ensino que fundamentem a observação realizada na escola. (REFERÊNCIAS)}
        
        \textcolor{red}{Caracterizar resumidamente o local de estágio, as pessoas envolvidas e as circunstâncias gerais que cercam o desenvolvimento do estágio de observação. Deve também conter a importância do trabalho (relatório), sua finalidade do estágio (objetivos), os procedimentos adotados (observação direta, consulta a documentos, conversas informais, participação em aulas em conjunto com o supervisor).}


        \textcolor{red}{Finalizar a introdução informando o que será apresentado nos capítulos seguintes.}

\newpage


%_______________A ESCOLA CAMPO DE ATUAÇÃO DOCENTE____________%
\section{A ESCOLA CAMPO DE ATUAÇÃO DOCENTE}

        \textcolor{red}{Neste item, o acadêmico deverá relatar os dados da observação, fundamentando teoricamente.}
        
        \textcolor{red}{Como ele percebe a escola campo de Estágio, questões gerais de sua organização física (no sentido de arquitetura), localização (periferia, centro, bairro), relação com o entorno, história, pessoas que fazem parte da comunidade escolar.}

    \subsection{Contextualização histórica da Unidade de Ensino}




    \subsection{Condições físicas e materiais da Unidade de Ensino}



    \subsection{Características básicas corpo técnico, administrativo e docente}

        \textcolor{red}{Número de profissionais, contratados e efetivos.}

    \subsection{Os estudantes atendidos pela escola}

        \textcolor{red}{Número de alunos, séries.}

    \subsection{A dinâmica do cotidiano escolar}

        \textcolor{red}{Turno de funcionamento. Como as pessoas da comunidade escolar se relacionam no dia a dia.}
        
        \textcolor{red}{Como as atividades escolares acontecem no cotidiano da escola.}
        
        \textcolor{red}{Horários de aula, intervalos.}

    \subsection{Projeto Político Pedagógico}

        \textcolor{red}{Alguns dos dados acima são do Projeto da Escola. Aqui podem ser colocadas outras questões, como filosofia da escola, objetivos, como está prevista a avaliação da aprendizagem, acordos, normatização sobre disciplina dos alunos, etc.}

\newpage

%_________________O ENSINO DE FÍSICA NA ESCOLA_________________%
\section{O ENSINO DE FÍSICA NA ESCOLA}

        \textcolor{red}{Neste capítulo o foco será o ensino de Física.}

    \subsection{Características do professor}

        \textcolor{red}{Poderão ser relatados aspectos da formação do professor, em quantas escolas trabalha, quantas turmas atende.}

        \textcolor{red}{Diante das condições que a escola oferece como ele realiza a docência.}

        \textcolor{red}{Se o professor segue o livro didático, ou usa sequência alternativa.}

    \subsection{Características das turmas}

        \textcolor{red}{Poderão ser especificados a faixa etária dos alunos da turma (ou turmas), série, atividades profissionais (no caso de estudantes trabalhadores), aspectos pedagógicos, etc.}

    \subsection{Interações professor-aluno}

        \textcolor{red}{Se houver oportunidade pode-se observar (e relatar) como os alunos de uma mesma turma se portam em diferentes disciplinas.}

    \subsection{Conteúdos curriculares de Física}

        \textcolor{red}{Quais conteúdos estão previstos. Aqui é interessante observar os Planos de Ensino para as três séries e analisar se determinados conteúdos de Física são deixados de fora ou se novos são incluídos.}

    \subsection{Os espaços de ensino-aprendizagem}

        \textcolor{red}{Os locais onde acontecem as aulas de Física. As metodologias utilizadas nas aulas, como é realizada a avaliação.}


\newpage

%___________A EXPERIÊNCIA DA DOCÊNCIA COMPARTILHADA____________%
\section{A EXPERIÊNCIA DA DOCÊNCIA COMPARTILHADA}

        \textcolor{red}{Escrever aqui como foi a inserção nas aulas de Física. A experiência de fazer atividades docentes junto com o supervisor.}

    \subsection{O planejamento}

        \textcolor{red}{Relatar o planejamento de atividades com o professor supervisor. Como foi a negociação. Como as ideias surgiram.}

    \subsection{Reflexão sobre a docência compartilhada}

        \textcolor{red}{A importância para a formação de participar de atividades com o professor supervisor. Análise da experiência.}


\newpage

%_______________CONSIDERAÇÕES FINAIS______________%
\section{CONSIDERAÇÕES FINAIS}

        \textcolor{red}{O acadêmico poderá escrever os problemas surgidos no estágio, limitações, pontos positivos, aspectos da formação que contribuíram para a atuação no campo profissional e as expectativas com relação aos próximos estágios.}

\newpage

\nocite{BNCC}

\bibliography{bibliografia}  %Preencha o arquivo bibliografia.bib com as referência adequadamente formatadas (pesquise sobre arquivos de extensão .bib) e descomente esta linha para que aparecem já em formato ABNT

\end{document}


\end{document}